\documentclass[10pt,a4paper]{amsart}
\usepackage{amsmath,amssymb,amsthm,graphicx}
% Retained preamble for the Stacks Project chapter {CH} (shared project preamble).
% The chapter file pinned-original-source.tex only carries \input{{preamble}} on its first line,
% so no preamble text lives inside the chapter itself.  The Stacks Project uses one single shared
% preamble for every chapter: the snapshot stages it three times (sources/stacks-chow/preamble.tex,
% sources/stacks-etale-cohomology/preamble.tex and focused-seeds/AGP/sources/stacks-snapshot-20261011/
% preamble.tex) and those copies are identical in size.  The block below is copied mechanically,
% line for line, from sources/stacks-chow/preamble.tex lines 14-18 and 154-217: the Xy-pic setup and
% the native theorem declarations together with the mathematical macros.  Nothing was added, removed
% or redefined; the wrapper class and the amsmath/amssymb/amsthm/graphicx base packages are supplied
% by the assembler.
% For commutative diagrams we use Xy-pic
% Chapter: stacks-spaces-more-morphisms.tex -- More on Morphisms of Spaces.
\usepackage[all]{xy}

% We use 2cell for 2-commutative diagrams.
\xyoption{2cell}

% Theorem environments.
%
\theoremstyle{plain}
\newtheorem{theorem}[subsection]{Theorem}
\newtheorem{proposition}[subsection]{Proposition}
\newtheorem{lemma}[subsection]{Lemma}

\theoremstyle{definition}
\newtheorem{definition}[subsection]{Definition}
\newtheorem{example}[subsection]{Example}
\newtheorem{exercise}[subsection]{Exercise}
\newtheorem{situation}[subsection]{Situation}

\theoremstyle{remark}
\newtheorem{remark}[subsection]{Remark}
\newtheorem{remarks}[subsection]{Remarks}

\numberwithin{equation}{subsection}

% Macros
%
\def\lim{\mathop{\mathrm{lim}}\nolimits}
\def\colim{\mathop{\mathrm{colim}}\nolimits}
\def\Spec{\mathop{\mathrm{Spec}}}
\def\Hom{\mathop{\mathrm{Hom}}\nolimits}
\def\Ext{\mathop{\mathrm{Ext}}\nolimits}
\def\SheafHom{\mathop{\mathcal{H}\!\mathit{om}}\nolimits}
\def\SheafExt{\mathop{\mathcal{E}\!\mathit{xt}}\nolimits}
\def\Sch{\mathit{Sch}}
\def\Mor{\mathop{\mathrm{Mor}}\nolimits}
\def\Ob{\mathop{\mathrm{Ob}}\nolimits}
\def\Sh{\mathop{\mathit{Sh}}\nolimits}
\def\NL{\mathop{N\!L}\nolimits}
\def\CH{\mathop{\mathrm{CH}}\nolimits}
\def\proetale{{pro\text{-}\acute{e}tale}}
\def\etale{{\acute{e}tale}}
\def\QCoh{\mathit{QCoh}}
\def\Ker{\mathop{\mathrm{Ker}}}
\def\Im{\mathop{\mathrm{Im}}}
\def\Coker{\mathop{\mathrm{Coker}}}
\def\Coim{\mathop{\mathrm{Coim}}}

% Boxtimes
%
\DeclareMathSymbol{\boxtimes}{\mathbin}{AMSa}{"02}

%
% Macros for moduli stacks/spaces
%
\def\QCohstack{\mathcal{QC}\!\mathit{oh}}
\def\Cohstack{\mathcal{C}\!\mathit{oh}}
\def\Spacesstack{\mathcal{S}\!\mathit{paces}}
\def\Quotfunctor{\mathrm{Quot}}
\def\Hilbfunctor{\mathrm{Hilb}}
\def\Curvesstack{\mathcal{C}\!\mathit{urves}}
\def\Polarizedstack{\mathcal{P}\!\mathit{olarized}}
\def\Complexesstack{\mathcal{C}\!\mathit{omplexes}}
% \Pic is the operator that assigns to X its picard group, usage \Pic(X)
% \Picardstack_{X/B} denotes the Picard stack of X over B
% \Picardfunctor_{X/B} denotes the Picard functor of X over B
\def\Pic{\mathop{\mathrm{Pic}}\nolimits}
\def\Picardstack{\mathcal{P}\!\mathit{ic}}
\def\Picardfunctor{\mathrm{Pic}}
\def\Deformationcategory{\mathcal{D}\!\mathit{ef}}
\usepackage{cleveref}
\crefname{originalreference}{original reference}{original references}
\Crefname{originalreference}{Original reference}{Original references}
\makeatletter
\newlabel{coherent-theorem-algebraization}{{089A}{source}{}{sourcepointer-1}{}}
\newlabel{coherent-theorem-algebraization@cref}{{[originalreference][0][]089A}{[1][0][]source}{}{}{}}
\makeatother
\crefname{originalreference}{original reference}{original references}
\Crefname{originalreference}{Original reference}{Original references}
\makeatletter
\newlabel{cohomology-lemma-h1-invertible}{{09NU}{source}{}{sourcepointer-2}{}}
\newlabel{cohomology-lemma-h1-invertible@cref}{{[originalreference][0][]09NU}{[1][0][]source}{}{}{}}
\makeatother
\crefname{originalreference}{original reference}{original references}
\Crefname{originalreference}{Original reference}{Original references}
\makeatletter
\newlabel{cohomology-proposition-vanishing-Noetherian}{{02UZ}{source}{}{sourcepointer-3}{}}
\newlabel{cohomology-proposition-vanishing-Noetherian@cref}{{[originalreference][0][]02UZ}{[1][0][]source}{}{}{}}
\makeatother
\crefname{originallemma}{original lemma}{original lemmas}
\Crefname{originallemma}{Original lemma}{Original lemmas}
\makeatletter
\newlabel{lemma-thickening-scheme}{{05ZR}{source}{}{sourcepointer-4}{}}
\newlabel{lemma-thickening-scheme@cref}{{[originallemma][0][]05ZR}{[1][0][]source}{}{}{}}
\makeatother
\crefname{originalreference}{original reference}{original references}
\Crefname{originalreference}{Original reference}{Original references}
\makeatletter
\newlabel{spaces-over-fields-lemma-codim-1-point-in-schematic-locus}{{0ADD}{source}{}{sourcepointer-5}{}}
\newlabel{spaces-over-fields-lemma-codim-1-point-in-schematic-locus@cref}{{[originalreference][0][]0ADD}{[1][0][]source}{}{}{}}
\makeatother
\crefname{originalreference}{original reference}{original references}
\Crefname{originalreference}{Original reference}{Original references}
\makeatletter
\newlabel{varieties-lemma-dim-1-proper-projective}{{0A26}{source}{}{sourcepointer-6}{}}
\newlabel{varieties-lemma-dim-1-proper-projective@cref}{{[originalreference][0][]0A26}{[1][0][]source}{}{}{}}
\makeatother
\title{Algebraization of a formal algebraic space of relative dimension at most one}
\author{The Stacks Project Authors}
\date{stacks-project:spaces-more-morphisms.tex:local-snapshot-20261011; GNU FDL 1.2 or later}
\begin{document}
\maketitle
\noindent\textit{Editorial note.} Counted span is the author's own proof environment, source lines 10355-10393, immediately following the statement at 10338-10353.\par\medskip
\noindent\textit{Editorial note.} Five cross-chapter pointers (Cohomology of Schemes, Cohomology, Varieties, Spaces over Fields) plus one same-chapter lemma are recorded in assembly.external\_references.\par\medskip
\noindent\textit{Editorial note.} The obstruction group and its vanishing are displayed inside the span.\par\medskip
\noindent\textit{Allowed standard prerequisites (editorial).}\begin{itemize}
\item Obstruction theory for invertible sheaves on thickenings. Cohomology, lemma cohomology-lemma-h1-invertible, cited at proof line 10380. Turns the lifting of an invertible sheaf along the thickening into the vanishing of a class in H\textasciicircum{}2 of the kernel.
\item Vanishing of the obstruction group in dimension at most one. Cohomology, proposition cohomology-proposition-vanishing-Noetherian, cited at proof lines 10385-10386. Kills the obstruction because X\_1 is a Noetherian scheme of dimension at most 1.
\item Grothendieck's algebraization theorem. Cohomology of Schemes, theorem coherent-theorem-algebraization, cited at proof lines 10389-10391. Produces the projective morphism of schemes X -> S with the compatible system of isomorphisms X\_n = S\_n x\_S X.
\item Realising the formal thickening by a scheme. chapter lemma lemma-thickening-scheme, cited in the proof. Lets the formal algebraic spaces X\_n be treated as thickenings of the scheme X\_1.
\item Proper projective and dimension-one facts about X\_1. Varieties, lemma varieties-lemma-dim-1-proper-projective and Spaces over Fields, lemma spaces-over-fields-lemma-codim-1-point-in-schematic-locus, cited in the proof. Control the geometry of the one-dimensional fibre so that the cohomological vanishing applies.
\end{itemize}\medskip
\ifdefined\hypertarget\hypertarget{sourcepointer-1}{}\fi
\noindent\textit{Source reference.} Original reference 089A: focused-seeds/geometry/S9-spaces-formal-space-proper-reldim-1/ref-tags.txt, source lines 5--5. Navigation pointer only: the counted proof cites this statement by \textbackslash{}ref as a standing input; the deduction the counted proof has to make is carried out inside the counted span, and no part of that deduction is delegated to this pointer.\par\medskip
\ifdefined\hypertarget\hypertarget{sourcepointer-2}{}\fi
\noindent\textit{Source reference.} Original reference 09NU: focused-seeds/geometry/S9-spaces-formal-space-proper-reldim-1/ref-tags.txt, source lines 6--6. Navigation pointer only: the counted proof cites this statement by \textbackslash{}ref as a standing input; the deduction the counted proof has to make is carried out inside the counted span, and no part of that deduction is delegated to this pointer.\par\medskip
\ifdefined\hypertarget\hypertarget{sourcepointer-3}{}\fi
\noindent\textit{Source reference.} Original reference 02UZ: focused-seeds/geometry/S9-spaces-formal-space-proper-reldim-1/ref-tags.txt, source lines 7--7. Navigation pointer only: the counted proof cites this statement by \textbackslash{}ref as a standing input; the deduction the counted proof has to make is carried out inside the counted span, and no part of that deduction is delegated to this pointer.\par\medskip
\ifdefined\hypertarget\hypertarget{sourcepointer-4}{}\fi
\noindent\textit{Source reference.} Original lemma 05ZR: sources/stacks-spaces-more-morphisms/pinned-original-source.tex, source lines 1730--1730. Navigation pointer only: the counted proof cites this statement by \textbackslash{}ref as a standing input; the deduction the counted proof has to make is carried out inside the counted span, and no part of that deduction is delegated to this pointer.\par\medskip
\ifdefined\hypertarget\hypertarget{sourcepointer-5}{}\fi
\noindent\textit{Source reference.} Original reference 0ADD: focused-seeds/geometry/S9-spaces-formal-space-proper-reldim-1/ref-tags.txt, source lines 8--8. Navigation pointer only: the counted proof cites this statement by \textbackslash{}ref as a standing input; the deduction the counted proof has to make is carried out inside the counted span, and no part of that deduction is delegated to this pointer.\par\medskip
\ifdefined\hypertarget\hypertarget{sourcepointer-6}{}\fi
\noindent\textit{Source reference.} Original reference 0A26: focused-seeds/geometry/S9-spaces-formal-space-proper-reldim-1/ref-tags.txt, source lines 9--9. Navigation pointer only: the counted proof cites this statement by \textbackslash{}ref as a standing input; the deduction the counted proof has to make is carried out inside the counted span, and no part of that deduction is delegated to this pointer.\par\medskip
\begin{lemma}
\label{lemma-formal-algebraic-space-proper-reldim-1}
Let $(A, \mathfrak m, \kappa)$ be a complete local Noetherian ring.
Set $S = \Spec(A)$ and $S_n = \Spec(A/\mathfrak m^n)$.
Consider a commutative diagram
$$
\xymatrix{
X_1 \ar[r]_{i_1} \ar[d] & X_2 \ar[r]_{i_2} \ar[d] & X_3 \ar[r] \ar[d] &
\ldots \\
S_1 \ar[r] & S_2 \ar[r] & S_3 \ar[r] & \ldots
}
$$
of algebraic spaces with cartesian squares. If $\dim(X_1) \leq 1$,
then there exists a projective morphism of schemes $X \to S$
and isomorphisms $X_n \cong X \times_S S_n$ compatible with $i_n$.
\end{lemma}

\begin{proof}
By Spaces over Fields, Lemma
\ref{spaces-over-fields-lemma-codim-1-point-in-schematic-locus}
the algebraic space $X_1$ is a scheme. Hence $X_1$
is a proper scheme of dimension $\leq 1$ over $\kappa$.
By Varieties, Lemma \ref{varieties-lemma-dim-1-proper-projective}
we see that $X_1$ is H-projective over $\kappa$.
Let $\mathcal{L}_1$ be an ample invertible sheaf on $X_1$.

\medskip\noindent
We are going to show that $\mathcal{L}_1$ lifts to a compatible system
$\{\mathcal{L}_n\}$ of invertible sheaves on $\{X_n\}$.
Observe that $X_n$ is a scheme too by Lemma \ref{lemma-thickening-scheme}.
Recall that $X_1 \to X_n$ induces homeomorphisms of underlying
topological spaces. In the rest of the proof we do not distinguish
between sheaves on $X_n$ and sheaves on $X_1$.
Suppose, given a lift $\mathcal{L}_n$ to $X_n$. We consider
the exact sequence
$$
1 \to
(1 + \mathfrak m^n\mathcal{O}_{X_{n + 1}})^* \to
\mathcal{O}_{X_{n + 1}}^* \to \mathcal{O}_{X_n}^* \to 1
$$
of sheaves on $X_{n + 1}$. The class of $\mathcal{L}_n$ in
$H^1(X_n, \mathcal{O}_{X_n}^*)$ (see
Cohomology, Lemma \ref{cohomology-lemma-h1-invertible})
can be lifted to an element of $H^1(X_{n + 1}, \mathcal{O}_{X_{n + 1}}^*)$
if and only if the obstruction in
$H^2(X_{n + 1}, (1 + \mathfrak m^n\mathcal{O}_{X_{n + 1}})^*)$
is zero. As $X_1$ is a Noetherian scheme of dimension $\leq 1$
this cohomology group vanishes (Cohomology, Proposition
\ref{cohomology-proposition-vanishing-Noetherian}).

\medskip\noindent
By Grothendieck's algebraization theorem
(Cohomology of Schemes, Theorem \ref{coherent-theorem-algebraization})
we find a projective morphism of schemes $X \to S = \Spec(A)$
and a compatible system of isomorphisms $X_n = S_n \times_S X$.
\end{proof}

\end{document}
